% Candidate: actual irreducible-residual point fields; no promotion yet.
\subsection{Point-field degrees and ramification in the irreducible residual case}

Let $p\ge5$, $q=p^2$, $N=q-1$ and $e=p(p-1)$. Let $K$ be
complete discretely valued of mixed characteristic $(0,p)$ with
algebraically closed residue field, $e_K=e$ and $\mu_p\subset K$.
Suppose that two whole potentially supersingular height-two groups
$M_1,M_2$ have an actual $p^3$-torsion isomorphism, irreducible
residual inertia, and different finite special characters. Let
$L/K$ be the common tame good extension of degree $N$, and put
\[
 F_2=L(M_1[p^2])=L(M_2[p^2]),\quad
 F_3=L(M_1[p^3])=L(M_2[p^3]).                         \tag{B-RD97.1}
\]
B79 gives $M_i[p](L)=M_i[p]$, the exact coordinate valuations
$e/q^{j-1}$ for points of order $p^j$, and the actual matrix
layers $W_1\subset W_2\subset M_2(\mathbf F_p)$. Write
$r_i=\dim W_i$. Every lower valuation is normalized on its
named field; a conductor is the last upper break, without adding
one. All extensions below are totally ramified.

\begin{lemma}[Every actual second- and third-point orbit is full]
\label{lem:B-RD97-degrees}
\status{Candidate; complete actual-field proof awaiting audits.}
For every primitive $p^2$-point $P$, every lift $P^+$ of order
$p^3$ with $pP^+=P$, and the corresponding points in both models,
the common actual point fields
\[
 E=L(P),\qquad T=L(P^+)
\]
satisfy
\[
 [E:L]=q,\qquad[T:E]=q,\qquad[T:L]=q^2.              \tag{B-RD97.2}
\]
The extensions $E/L$ and $T/E$ are Galois, each with all $q$
first-level translations. In particular
\[
       W_1v=\mathbf F_p^2\quad\text{for every }v\ne0,
       \qquad r_1\ge2.                              \tag{B-RD97.3}
\]
No assertion that $T/L$ is Galois is made here.
\end{lemma}
\begin{proof}
Since $pP$ is $L$-rational and all first-level translations are
$L$-rational, every conjugate of $P$ lies in $L(P)$. Thus $E/L$
is Galois with translation group $W_1v$, where $v=pP$ is viewed
in the first-level module. The valuation $v_L(P)=(p-1)/p$
forces $[E:L]\ge p$, and the translation group has order at most
$q$. Likewise $T/E$ is Galois with a subgroup of the same
$L$-rational first-level module as its translations, so its degree
is at most $q$.

Suppose $[E:L]=p$. The valuation
$v_L(P^+)=(p-1)/p^3$ forces $[T:L]\ge p^3$. The preceding
upper bounds give equality, and $\operatorname{Gal}(T/E)$ has
all $q$ translations. In $T$, the two coordinate functions $x_i$
of $P^+$ have valuation $A=p-1$, prime to $p$. A nonidentity
translation displaces each by a nonzero first-level coordinate,
of valuation $D=ep^3$. The fixed-coefficient equality of B60
therefore gives one common lower break $D-A$ and
\[
       \overline{(gx_i-x_i)/\varpi^D}
                     =\overline A\,a_i d_g,
       \qquad a_i\ne0.                              \tag{B-RD97.4}
\]
Here $\varpi$ is one uniformizer of the actual common field $T$,
$a_i=\overline{x_i/\varpi^A}$, and
$d_g=\overline{(g\varpi-\varpi)/\varpi^{D-A+1}}$.
The fixed residue representatives can be taken in the base $E$;
all displacements are below $v_T(p)=ND$.

We recall precisely the leading-space comparison. In Honda
coordinates the first-level leading spaces are scalar multiples
of $\mathbf F_q$. Their tame exponents are $c_i+e$, as in B79.
An actual residual isomorphism between unequal special characters
has Frobenius orientation on these spaces: after nonzero scalar
normalizations its parameters satisfy $z_2=z_1^p$. A scalar
orientation would instead give $c_1=c_2\pmod N$. The statement
is unchanged on replacing a base uniformizer by a uniformizer of
a point field: this multiplies each leading space by a fixed
nonzero scalar. Equation (B-RD97.4) says that the two leading
translation maps are scalar multiples of the same $d_g$.
Their values fill the whole first-level space, since every
translation occurs. Frobenius on $\mathbf F_q$ would be a fixed
scalar map, contradicted by $1$ and an element outside $\mathbf F_p$.
Thus $[E:L]=p$ is impossible, proving $[E:L]=q$ and (B-RD97.3).

The exact identity
$(1+pA+p^2B)^p=1+p^2A\pmod{p^3}$ gives $W_1\subset W_2$
in the same matrix coordinates. Hence $W_2v$ already equals the
full first-level module. Terminal matrices fix $P$ and translate
$P^+$ by every first-level point. Thus $T/E$ has all $q$
translations and degree $q$, proving the remaining assertions.
This is an actual point-field argument, not an assumed sum of
residual evaluation dimensions.
\end{proof}

\begin{proposition}[Two first-layer dimensions are impossible]
\label{prop:B-RD97-two}
\status{Candidate; complete actual-field proof awaiting audits.}
Under the standing comparison hypotheses, $r_1\ge3$.
\end{proposition}
\begin{proof}
Suppose $r_1=2$ and fix a primitive point $P^+$. By the preceding
lemma, evaluation $W_1\longrightarrow\mathbf F_p^2$ at its
nonzero residual vector is an isomorphism. Every element of the
actual point stabilizer in $\operatorname{Gal}(F_3/L)$ therefore
lies in the terminal subgroup $1+p^2W_2$. That subgroup is central
in the matrix kernel modulo $p^3$. The point stabilizer is thus
normal, so $T=L(P^+)/L$ is Galois. Moreover $E=F_2$ and
\[
 |G|=p^4,\quad G=\operatorname{Gal}(T/L),\quad
 U=\operatorname{Gal}(T/E)\simeq\mathbf F_p^2,
 \quad G/U\simeq\mathbf F_p^2.                      \tag{B-RD97.5}
\]
The actual $p$th-power map onto $U$ is surjective: it is evaluation
of $W_1$ under the exact matrix identity just used. The group $G$
need not be abelian, and no Cartan structure is asserted or used.

In $T$ the primitive, second-level and first-level coordinate
valuations are respectively
\[
 A=e=p(p-1),\qquad D_1=qA,\qquad D_2=q^2A,
                      v_T(p)=ND_2.                  \tag{B-RD97.6}
\]
Every $g\notin U$ displaces the selected point by a point of exact
order $p^2$, since its first-layer evaluation is nonzero. Every
$h\in U\setminus\{1\}$ displaces it by a nonzero first-level
point. Normalize the leading coefficient of the selected point
as $a_i\ne0$. The leading $p$-series term is uniquely $T^q$.
Normalize the first-level leading space using the actual common
point $p^2P^+$ as parameter1. Then the normalized second-level
space uses $pP^+$ as parameter1. The leading displacements are
\[
 a_i^q\lambda_i(g),\qquad
 a_i^{q^2}\lambda_i(h),\qquad
                  \lambda_2=\lambda_1^p.             \tag{B-RD97.7}
\]
The first parameter map is an additive isomorphism on $G/U$,
and the second is one on $U$. The Frobenius identity follows
from the leading first-level orientation in the preceding proof;
taking unique $q$th roots transports it to the second level.
This uses only one selected primitive point in each model.

Fix their full coordinate expansions with fixed residue
representatives, and let $a$ be the smallest lower break in $G$.
Every nonidentity element of $U$ has break greater than $a$:
each is a $p$th power, and the additive first ramification
coefficient vanishes on a $p$th power, so $b(g^p)>b(g)$.
Put
\[
 a_0=(D_1-A)/p=N(p-1),\qquad
 j_0=A+(p-1)a_0=D_1-a_0,\qquad j_0\equiv-1\pmod p.
                                                        \tag{B-RD97.8}
\]
Because $v_p(A)=1$, at a break $b$ the first possible coordinate
change is $\min\{A+pb,J_i+b\}$, where $J_i$ is its first
nonzero ordinary exponent, possibly infinite. All calculations
through $D_2$ are below $v_T(p)$.

For a minimum-break $g\notin U$, if $a<a_0$, cancellation below
$D_1$ forces $J_i=A+(p-1)a$ with a nonzero coefficient in both
models. If $a>a_0$, the actual change at $D_1$ forces
$J_i=D_1-a$, again with nonzero coefficient. An impermissible
or $p$-divisible asserted first ordinary degree already gives a
contradiction. At every terminal break $b>a$, the fixed ordinary
term is uniquely earlier than every divisible term: the differences
are $(p-1)(b-a)$ in the first case and
$(p-1)(b-a)+p(a-a_0)$ in the second. It must therefore give
$J_i+b=D_2$. Its leading coefficient gives
$\lambda_i(h)=\kappa_i d_h$ with $\kappa_i\ne0$ independent
of $h$. Since the values of $\lambda_1$ on $U$ fill $\mathbf F_q$,
this contradicts (B-RD97.7). Thus $a=a_0$.

At this boundary no ordinary degree below $j_0$ can occur.
Write $f_i$ for its possibly zero coefficient. At degree $D_1$
the only two possible terms give, for a fixed minimum element,
\[
 a_i^q\lambda_i(g)=(A/p)a_i d_g^p+j_0 f_i d_g.
                                                        \tag{B-RD97.9}
\]
If one $f_i$ is nonzero, every terminal break is
$b_*=D_2-j_0>a_0$ and that ordinary term is uniquely first.
In the other model the same break forces that exact nonzero
ordinary term as well, since
\[
           b_*-(D_2-A)/p=(p-1)(D_2-D_1)/p>0.
                                                        \tag{B-RD97.10}
\]
The same scalar-versus-Frobenius contradiction on $U$ follows.
Hence $f_1=f_2=0$. No nonterminal element can then have break
greater than $a_0$: all its divisible and ordinary changes would
be after $D_1$. The nonterminal elements all have break $a_0$
and the terminal ones have greater break, so the leading
ramification map identifies $G/U$ with a two-dimensional additive
subspace of the residue field. For all its elements (B-RD97.9)
reduces to $\lambda_i(g)=(A/p)a_i^{1-q}d_g^p$.
The two maps are scalar proportional, contrary to their Frobenius
relation and full $\mathbf F_q$ image. This proves the proposition.
It is the complete two-level fixed-coefficient argument of B81
on an actual single-point Galois field; no scalar dependence on a
family of primitive lifts, abelian image, or endomorphism is needed.
\end{proof}

\begin{lemma}[Exact last conductors of the two actual point towers]
\label{lem:B-RD97-conductors}
\status{Candidate; complete actual-field proof awaiting audits.}
For every actual tower in (B-RD97.2), the last upper breaks are
\[
 c_1=N(p-1)\quad\text{for }E/L,\qquad
 c_2=(q^2-1)(p-1)\quad\text{for }T/E.
                                                        \tag{B-RD97.11}
\]
The full second-level extension $F_2/L$ also has last upper
break $c_1$. The last upper break of $F_3/F_2$ is at most
\[
                         p^{r_1-2}c_2.               \tag{B-RD97.12}
\]
\end{lemma}
\begin{proof}
In either point extension its group is $\mathbf F_p^2$, its
coordinate valuation is $A=e$ with $v_p(A)=1$, and every
nonidentity element has the same coordinate displacement $D$.
Here $D=qA$ for $E/L$ and $D=q^2A$ for $T/E$. In both cases
$v(p)=ND>D$. The actual leading translation maps in the two
models are related by Frobenius and fill the entire normalized
first-level space. Put $b_0=(D-A)/p$.

With one lower break $b<b_0$, the first-change cancellation
forces $J=A+(p-1)b$ and a relation $d_g^{p-1}=\kappa\ne0$
on all nonzero leading coefficients of the group. The latter
form a two-dimensional additive space, so there are $q-1$
values, more than the $p-1$ roots of this polynomial.
If $b>b_0$, a unique ordinary term at degree $D-b$ makes the
two leading translation maps scalar multiples of the same $d_g$,
contrary to their Frobenius relation. Thus a single break is $b_0$.

With two lower breaks $t<u$, one must have $t<b_0$. Indeed if
$t\ge b_0$, no ordinary term can change before $D$ at $t$
(it would be uniquely first), and every term then changes after
$D$ at $u$, impossible. Cancellation at $t<b_0$ forces the
first ordinary degree $J=A+(p-1)t$. It is uniquely first at $u$,
so
\[
                  J+u=D,\qquad u+(p-1)t=D-A.
                                                        \tag{B-RD97.13}
\]
The largest cyclic quotient break is $(u+(p-1)t)/p=b_0$ by
the exact quotient average. This proves (B-RD97.11), including
the one-break boundary and all absent-coefficient cases.

Choose $P_1,P_2$ forming a basis of the actual $p^2$-module and
write $E_i=L(P_i)$. Then $F_2=E_1E_2$ as actual fields.
Each is Galois elementary abelian of degree $q$ with last
conductor $c_1$. Every cyclic character of their compositum is
a sum of characters of the two fields, whose conductor is at
most their maximum. Thus the largest conductor of $F_2/L$ is
at most $c_1$, and is at least $c_1$ because it has either
$E_i$ as quotient. This proves its exact last break.

Now lift these to a basis $P_i^+$ of the actual $p^3$-module,
and put $T_i=L(P_i^+)$ and $D_i=F_2T_i$. The terminal matrices
already give all translations of $T_i/E_i$, so
\[
 F_2\cap T_i=E_i,\qquad [D_i:F_2]=q,
 \qquad[F_2:E_i]=p^{r_1-2},\qquad D_1D_2=F_3.
                                                        \tag{B-RD97.14}
\]
For clarity, if $H_i$ is the stabilizer of $P_i$ and $S_i$
that of $P_i^+$ in the full actual group, then $US_i=H_i$,
where $U$ is the terminal subgroup. This proves the intersection
and degree assertions. The subgroup $U\cap S_i$ is central
and normal in $H_i$. For first-layer matrices $A,C$ in $H_i$,
$Av_i=Cv_i=0$, whence $[A,C]v_i=0$. The exact commutator and
$p$th-power formulas put both in $U\cap S_i$. Thus $D_i/E_i$
is an actual elementary-abelian compositum; no abstract disjoint
extensions have been substituted.

Herbrand subgroup conversion bounds every conductor of
$D_i/F_2$ by $\psi_{F_2/E_i}(c_2)$. One direct justification
is to take upper index $u>c_2$ in $\operatorname{Gal}(D_i/E_i)$.
Its image in $\operatorname{Gal}(T_i/E_i)$ is trivial; its
intersection with $\operatorname{Gal}(D_i/F_2)$ is therefore
trivial because $D_i=F_2T_i$. The subgroup conversion now gives
the asserted bound. The inverse Herbrand function has slopes
at most $[F_2:E_i]=p^{r_1-2}$ and starts at zero, so this is
at most $p^{r_1-2}c_2$. Finally the two actual quotients
$D_i/F_2$ generate the full elementary-abelian terminal field
$F_3/F_2$; the same character-sum bound proves (B-RD97.12).
Only the actual subgroup/quotient formulas of
\cite[Proposition3.7, Theorem3.10 and Proposition3.13,
complete proofs, pp.~14--16]{ColmezCorpsLocaux} are used.
\end{proof}

\begin{proposition}[The first matrix layer is not traceless]
\label{prop:B-RD97-trace}
\status{Candidate; complete actual-field proof awaiting audits.}
Under the standing comparison hypotheses, the trace on $W_1$
is nonzero.
\end{proposition}
\begin{proof}
Otherwise $r_1\ge3$ and $W_1\subset\mathfrak{sl}_2$ give
$W_1=\mathfrak{sl}_2$ and $r_1=3$. The actual determinant on
the good field $L$ is the cyclotomic character. Tracelessness
of the first layer makes it trivial modulo $p^2$, so
$\mu_{p^2}\subset L$. The full arbitrary-residue proof of
Lemma~\ref{lem:B-RD86-cyclotomic} gives an actual cyclic degree-$p$
field $C=L(\mu_{p^3})$ with break $N^2$. That proof identifies
$K=K_0(\mu_{p^2})$ using a Cohen coefficient ring and the exact
absolute index; it is not a finite-residue local-field substitution.

Since $N^2>c_1$ and the full $F_2/L$ has last upper break $c_1$,
the fields $C$ and $F_2$ are disjoint over $L$. Their actual
elementary-abelian compositum has $C$ as its unique latest
character direction. Hence the restricted cyclotomic character
on $F_3/F_2$ has conductor exactly
\[
 \psi_{F_2/L}(N^2)
   =\psi_{F_2/L}(c_1)+p^3(N^2-c_1)
   >p^3(N^2-c_1).                                   \tag{B-RD97.15}
\]
It is a nonzero actual terminal character since $C\subset F_3$
and $C\not\subset F_2$. But the preceding lemma bounds every
terminal conductor by $p c_2$. The exact positive difference is
\[
 p^3(N^2-c_1)-p c_2
          =pN(p-1)(p^3-p^2-1)>0.                    \tag{B-RD97.16}
\]
This is the required contradiction. No unknown cyclotomic break
over an arbitrary base has been assigned: the exact cyclotomic
field here follows from the assumed traceless first layer.
\end{proof}

\begin{corollary}[Only traceful seven- or eight-dimensional images remain]
\label{cor:B-RD97-layers}
\status{Candidate; complete actual-field proof awaiting audits.}
An irreducible-residual different-character comparison as above
must have
\[
 (r_1,r_2)=(3,4)\text{ with }\operatorname{tr}|_{W_1}\ne0,
 \quad\text{or }(r_1,r_2)=(4,4).                    \tag{B-RD97.17}
\]
At $p=5,e_K=20$, together with B90 and B96 this leaves only
these irreducible-residual wild images of orders $5^7$ and $5^8$
for a whole potentially supersingular counterexample at level125.
\end{corollary}
\begin{proof}
We know $3\le r_1\le r_2\le4$. If $r_1=r_2=3$, the common
space $V=W_1=W_2$ is a Lie subalgebra by the exact matrix
commutator formula, normalized by an irreducible residual tame
matrix. We verify that it must be $\mathfrak{sl}_2$.
If its trace is nonzero, $V\cap\mathfrak{sl}_2$ is a
two-dimensional Lie algebra invariant under that tame matrix.
It cannot be abelian: the centralizer of any nonzero traceless
matrix has dimension two in $M_2$ and intersects $\mathfrak{sl}_2$
in a line, because $p\ne2$ and the centralizer contains the identity.
Its derived line therefore has a generator $Y$ and an element
$X$ with $[X,Y]=Y$. Taking traces of this equation and of
$[X,Y^2]=2Y^2$ gives $\operatorname{tr}Y=
\operatorname{tr}Y^2=0$, hence $\det Y=0$.
Thus $Y$ is nonzero nilpotent. Its line and its kernel are
preserved by the tame normalizer, contradicting residual
irreducibility. This proves $V=\mathfrak{sl}_2$, which was
excluded by the preceding proposition. The remaining layer
shapes are exactly (B-RD97.17). Their orders are the products
of the two actual layer orders, as in B79.
\end{proof}

\begin{remark}[Scope and intended supplement]
\label{rem:B-RD97-scope}
\status{Candidate; accompanies the unpromoted point-field results.}
The two remaining layer shapes are necessary possibilities only.
No actual formal-group comparison is constructed or excluded in
either remaining shape. At $p=5,e=20$ the unrestricted least index
still belongs to $\{2,3\}$. The fixed-point argument does not
replace a finite-level representation by a CM group or impose an
abelian image. Every field and relative translation group used
above is supplied by the original actual $p^3$-comparison.

The complete B79 leading-coordinate and orientation proof, B80
integral-lift and coefficient arguments, B81 two-level boundary
proof, B86 Cohen/cyclotomic lemma and B90 actual point-field
intersection proof were reread. No new external-source opening
is claimed. The separate standalone exact supplement
\path{code/reduction_descent/verify_irreducible_point_field_layers.py}
passed at 05:24:04.441182 UTC on 20 November 2026 in
1.2538034710014472 seconds. At $p=5,7,11$ it enumerates all
19,882 matrix planes, identifies 7,132 planes with every nonzero
evaluation surjective, and checks 934,848 corresponding nonzero
evaluations and 88,335 literal matrix powers. It checks 14,109
permitted first-break rows, 1,584 impermissible first degrees, all
nine displayed boundary comparisons and 164,874 uniform point
conductor rows. It enumerates 2,020 matrix hyperplanes and their
29 Lie subalgebras, verifying the common invariant line in every
nontraceless case. The quadratic-field checks use all 3,345
elements at nine stated primes, alongside the exact cyclotomic
inequalities. These finite checks supplement the complete all-prime
proof and infer no field or formal-group realization. The exact
executed source, complete output and stdout are retained; no
executed checker failure occurred.
\end{remark}
